%===================================================================
\section{Evaluation Strategy}
%===================================================================

\subsection{Evaluation Overview}

Evaluation of literature review systems is a notoriously difficult problem because there is no single ground truth review for a given topic. Our evaluation strategy addresses this through four complementary methods: automated metric computation against expert surveys, novel task-specific metrics, human expert evaluation, and the empirical gap validation protocol.

\subsection{Benchmark Datasets and Domains}

Three domains are selected representing different characteristics:

\begin{enumerate}
    \item \textbf{Retrieval-Augmented Generation (RAG):} A rapidly evolving domain with frequent contradictions and active methodological debate.
    \item \textbf{Graph Neural Networks (GNN):} A more mature domain with established consensus on core findings and well-defined evaluation benchmarks.
    \item \textbf{Low-Resource Machine Translation:} A domain with significant methodological diversity and cross-lingual challenges.
\end{enumerate}

For each domain, a corpus of 50--150 primary research papers spanning 2017--2025 is assembled via Semantic Scholar and arXiv. Gold standard expert-authored surveys are withheld from the retrieval corpus and used only for evaluation. A held-out set of papers published in 2025--2026 is assembled for gap validation.

\subsection{Novel Task-Specific Metrics}

We introduce five evaluation metrics designed to assess the specific capabilities of uncertainty-aware synthesis:

\begin{description}
    \item[Claim Coverage Score (CCS):] Measures the proportion of key claims identified in the gold standard survey that are present (in semantically equivalent form) in the generated review. Computed using embedding similarity with threshold matching. CCS measures content recall from an expert perspective.

    \item[Confidence Calibration Score (CalibScore):] Measures whether confidence assignments are calibrated. For claims marked as High Confidence, the proportion that experts agree are well-established should be high. Calibration curves are plotted for each confidence category to assess alignment between system scores and expert consensus.

    \item[Contradiction Detection F1 (CD-F1):] A held-out contradiction pair set is constructed by domain experts identifying genuine contradictions within each corpus. Precision is the proportion of system-identified contradictions confirmed by experts; recall is the proportion of expert-identified contradictions found by the system.

    \item[Gap Precision at K (GP@K):] The top-K gap objects generated by the system are evaluated against the held-out future paper set. A gap is considered validated if a paper in the held-out set directly addresses the gap as specified by the falsifiability statement. GP@5 and GP@10 are reported.

    \item[Temporal Coherence Score (TCS):] Expert reviewers assess whether the temporal narrative produced by the TKET accurately represents the historical development of the field, rated on a 5-point Likert scale.
\end{description}

\subsection{Human Evaluation Protocol}

Human evaluation is conducted with five domain experts per domain. Reviewers are presented with the generated review alongside the gold standard survey and baseline system outputs. Evaluation dimensions include:

\begin{itemize}
    \item \textbf{Factual Accuracy:} Are claims accurately attributed and correctly stated?
    \item \textbf{Synthesis Quality:} Does the review integrate multiple papers rather than merely summarizing individually?
    \item \textbf{Contradiction Handling:} Are conflicting findings acknowledged and addressed?
    \item \textbf{Gap Usefulness:} Are identified research gaps specific and actionable?
    \item \textbf{Temporal Awareness:} Does the review accurately convey the historical development of the field?
\end{itemize}

Each dimension is rated on a 7-point Likert scale. Inter-rater reliability is measured using Cohen's Kappa. Reviewers are blind to which system generated which output.

\subsection{Ablation Study Design}

To validate the contribution of each component, a systematic ablation study compares the following configurations (Table~\ref{tab:ablation}):

\begin{table}[ht]
\caption{Ablation study configurations. \checkmark\ indicates component inclusion.}\label{tab:ablation}
\centering
\small
\begin{tabular}{lccccc}
\toprule
\textbf{Configuration} & \textbf{CGCERS} & \textbf{CDRA} & \textbf{FRGO} & \textbf{DSKG} & \textbf{TKET} \\
\midrule
Full System & \checkmark & \checkmark & \checkmark & \checkmark & \checkmark \\
$-$CGCERS & & \checkmark & \checkmark & \checkmark & \checkmark \\
$-$CDRA & \checkmark & & \checkmark & \checkmark & \checkmark \\
$-$FRGO & \checkmark & \checkmark & & \checkmark & \checkmark \\
$-$DSKG (flat retrieval) & \checkmark & \checkmark & \checkmark & & \checkmark \\
$-$TKET & \checkmark & \checkmark & \checkmark & \checkmark & \\
Baseline (RAG-only) & & & & & \\
\bottomrule
\end{tabular}
\end{table}

%===================================================================
\section{Limitations}
%===================================================================

Despite the contributions of this work, several limitations should be acknowledged:

\begin{enumerate}
    \item \textbf{LLM Hallucination Risk in Claim Extraction:} Despite structured prompting, LLMs may hallucinate claim details not present in the source paper. This is mitigated through post-extraction verification that checks extracted claims against source text using substring and semantic matching, but cannot be fully eliminated.

    \item \textbf{Knowledge Graph Construction Accuracy:} Relationship classification errors during DSKG construction may introduce incorrect edges that propagate through synthesis. The Critic Agent's coverage evaluation partially mitigates this but cannot catch all errors.

    \item \textbf{Scalability:} The contradiction detection pairwise comparison has quadratic complexity $O(|\mathcal{C}|^2)$ in the number of claims. For corpora exceeding 500 papers, approximate nearest-neighbor search is used for candidate filtering, which may reduce recall.

    \item \textbf{Evaluation Ground Truth Limitations:} Expert-authored survey papers used as gold standards may themselves contain biases, omissions, and errors, limiting their validity as absolute ground truth.

    \item \textbf{Domain Generalization:} The pre-trained NLI model is primarily trained on NLP and computer science text. Performance on domains with different rhetorical conventions (e.g., medicine, social sciences) may require domain-specific fine-tuning.

    \item \textbf{PDF Processing Quality:} GROBID's section segmentation accuracy varies with PDF formatting quality, particularly for older papers without standardized layouts.
\end{enumerate}

%===================================================================
\section{Future Work}
%===================================================================

Several directions for future work are envisioned:

\begin{itemize}
    \item \textbf{Citation Intent Classification:} Incorporating citation intent classification (supportive, contrasting, neutral) from tools such as SciCite directly into DSKG edge construction for more accurate relationship modeling.

    \item \textbf{Domain-Specific Fine-Tuning:} Fine-tuning the Claim Extractor and Contradiction Detector agents on domain-specific corpora for medicine, law, and social sciences.

    \item \textbf{Real-Time Monitoring:} Extending the system to monitor newly published papers and automatically update the DSKG, Conflict Registry, and Gap Set when new work resolves a gap or introduces a new contradiction.

    \item \textbf{Peer Reviewer Simulation:} Adding a Peer Reviewer Agent that evaluates the generated review against venue-specific standards, simulating the review process for target publication venues.

    \item \textbf{Multi-Language Support:} Extending retrieval and processing to non-English language papers, important for fields with significant non-English publication traditions.

    \item \textbf{Interactive Refinement:} Enabling users to provide real-time feedback on specific claims, contradictions, or gaps through the Streamlit interface, creating a human-in-the-loop refinement cycle.
\end{itemize}

%===================================================================
\section{Conclusion}
%===================================================================

This paper has presented a novel autonomous multi-agent framework for scientific literature review generation that addresses three fundamental limitations of existing systems: the absence of epistemic grounding for synthesized claims, the inability to detect and resolve contradictions across papers, and the generation of vague, unverifiable research gap statements.

The system's five core technical contributions---Confidence-Graded Claim Extraction and Reliability Scoring (CGCERS), the Contradiction Detection and Resolution Agent (CDRA), the Formalized Research Gap Ontology (FRGO), the Dynamic Scientific Knowledge Graph (DSKG), and the Temporal Knowledge Evolution Tracker (TKET)---collectively advance the state of the art in automated scientific reasoning. The CGCERS formula (Equation~\ref{eq:cgcers}) provides calibrated confidence scores based on multi-paper consensus, temporal recency, and counter-evidence. The CDRA's two-stage architecture combining NLI model filtering with LLM confirmation enables efficient yet accurate contradiction detection with a structured four-dimensional taxonomy. The FRGO introduces empirically evaluable gap objects with explicit falsifiability statements---a first in the automated survey generation literature.

The proposed evaluation framework, introducing CCS, CD-F1, GP@K, TCS, and CalibScore, addresses the longstanding evaluation inadequacy in automated literature review research and provides a reusable assessment protocol for future work. The complete system has been implemented using Google Gemini as the primary LLM engine with a modular architecture supporting provider switching, NetworkX for graph-based reasoning, ChromaDB for semantic search, and GROBID for structured PDF processing, all orchestrated by an autonomy loop that iteratively refines outputs based on critic feedback.

By moving beyond static pipeline summarization toward a graph-centric, epistemically grounded, contradiction-aware, and temporally sensitive synthesis architecture, this work takes a meaningful step toward AI systems capable of genuine scientific reasoning partnership---not merely retrieving and rewording existing text, but reasoning about the structure of knowledge itself.

\subsubsection{Acknowledgements.} This work was conducted as part of the Generative AI course project at FAST-NUCES under the supervision of Dr.~Akhtar Jameel.

%===================================================================
% Bibliography
%===================================================================
\begin{thebibliography}{20}

\bibitem{ref_growth}
Johnson, R., Watkinson, A., Mabe, M.: The STM Report: An Overview of Scientific and Scholarly Publishing. 5th edn. International Association of Scientific, Technical and Medical Publishers (2018)

\bibitem{ref_autosurvey}
Li, Y., Zhang, Y., Sun, L.: AutoSurvey: Large Language Models Can Automatically Write Surveys. arXiv preprint arXiv:2406.10252 (2024)

\bibitem{ref_litllm}
Agarwal, S., Laforte, S., Bhatt, V.: LitLLM: A Toolkit for Scientific Literature Review. arXiv preprint arXiv:2402.01788 (2024)

\bibitem{ref_surveyagent}
Liu, Z., Lin, W., Shi, Y., Zhao, J.: SurveyAgent: A Conversational System for Personalized and Efficient Research Survey. arXiv preprint arXiv:2404.06364 (2024)

\bibitem{ref_researchagent}
Baek, J., Jauhar, S.K., Cucerzan, S., Hwang, S.J.: ResearchAgent: Iterative Research Idea Generation over Scientific Literature with Large Language Models. arXiv preprint arXiv:2404.07738 (2024)

\bibitem{ref_rag}
Lewis, P., Perez, E., Piktus, A., Petroni, F., Karpukhin, V., Goyal, N., K\"{u}ttler, H., Lewis, M., Yih, W., Rockt\"{a}schel, T., Riedel, S., Kiela, D.: Retrieval-Augmented Generation for Knowledge-Intensive NLP Tasks. In: Advances in Neural Information Processing Systems, vol.~33, pp. 9459--9474 (2020)

\bibitem{ref_rouge}
Lin, C.Y.: ROUGE: A Package for Automatic Evaluation of Summaries. In: Text Summarization Branches Out, pp. 74--81. ACL (2004)

\bibitem{ref_extractive}
Nenkova, A., McKeown, K.: A Survey of Text Summarization Techniques. In: Mining Text Data, pp. 43--76. Springer (2012). \doi{10.1007/978-1-4614-3223-4\_3}

\bibitem{ref_scisummnet}
Yasunaga, M., Kasai, J., Zhang, R., Fabbri, A.R., Li, I., Friedman, D., Radev, D.R.: ScisummNet: A Large Annotated Corpus and Content-Impact Models for Scientific Paper Summarization with Citation Networks. In: AAAI, pp. 7386--7393 (2019)

\bibitem{ref_react}
Yao, S., Zhao, J., Yu, D., Du, N., Shafran, I., Narasimhan, K., Cao, Y.: ReAct: Synergizing Reasoning and Acting in Language Models. In: ICLR (2023)

\bibitem{ref_reflexion}
Shinn, N., Cassano, F., Gopinath, A., Narasimhan, K., Yao, S.: Reflexion: Language Agents with Verbal Reinforcement Learning. In: NeurIPS, vol.~36 (2023)

\bibitem{ref_langgraph}
LangGraph: Build Stateful, Multi-Actor Applications with LLMs. \url{https://github.com/langchain-ai/langgraph} (2024). Last accessed 15 Apr 2026

\bibitem{ref_crewai}
CrewAI: Framework for Orchestrating Role-Playing Autonomous AI Agents. \url{https://github.com/crewAIInc/crewAI} (2024). Last accessed 15 Apr 2026

\bibitem{ref_orkg}
Jaradeh, M.Y., Oelen, A., Farfar, K.E., Prinz, M., D'Souza, J., Kismih\'{o}k, G., Stocker, M., Auer, S.: Open Research Knowledge Graph: Next Generation Infrastructure for Semantic Scholarly Knowledge. In: K-CAP, pp. 243--246. ACM (2019). \doi{10.1145/3360901.3364435}

\bibitem{ref_scholarlykg}
Dessì, D., Osborne, F., Reforgiato Recupero, D., Buscaldi, D., Motta, E.: AI-KG: An Automatically Generated Knowledge Graph of Artificial Intelligence. In: ISWC, LNCS, vol.~12507, pp. 127--143. Springer (2020). \doi{10.1007/978-3-030-62466-8\_9}

\bibitem{ref_gapsurvey}
Mohammad, S., Dorr, B., Egan, M., Hassan, A., Muthukrishan, P., Qazvinian, V., Radev, D., Zajic, D.: Using Citations to Generate Surveys of Scientific Paradigms. In: NAACL-HLT, pp. 584--592. ACL (2009)

\bibitem{ref_networkx}
Hagberg, A.A., Schult, D.A., Swart, P.J.: Exploring Network Structure, Dynamics, and Function Using NetworkX. In: Proceedings of the 7th Python in Science Conference, pp. 11--15 (2008)

\bibitem{ref_louvain}
Blondel, V.D., Guillaume, J.L., Lambiotte, R., Lefebvre, E.: Fast Unfolding of Communities in Large Networks. Journal of Statistical Mechanics: Theory and Experiment \textbf{2008}(10), P10008 (2008). \doi{10.1088/1742-5468/2008/10/P10008}

\bibitem{ref_deberta}
He, P., Liu, X., Gao, J., Chen, W.: DeBERTa: Decoding-Enhanced BERT with Disentangled Attention. In: ICLR (2021)

\bibitem{ref_gemini}
Gemini Team, Google: Gemini: A Family of Highly Capable Multimodal Models. arXiv preprint arXiv:2312.11805 (2024)

\end{thebibliography}

\end{document}
